\chapter{Deploying the Tier 2 Distilled Swarm: 1,000,000 Agents on Shared MLPs}
\label{ch:distilled_swarm_scaling}

The fundamental bottleneck of financial agent-based simulations has historically resided in the trade-off between cognitive fidelity and computational scale. Chapter \ref{ch:generative_swarms} deconstructed the physical ceiling that prevents online autoregressive Large Language Models (LLMs) from scaling beyond hundreds of concurrent instances. This chapter establishes the engineering foundation of Tier 2: deploying a distilled policy swarm of 1,000,000 concurrent agents executing on host CPU silicon at deterministic, sub-0.2 microsecond latency.

\section{Behavioral Archetype Factorization}

Simulating one million independent neural network parameter sets would overwhelm memory buses and prevent efficient instruction caching. In real-world financial markets, trading entities exhibit clustered microstructural behaviors dictated by mandates, risk constraints, and latency horizons. AMPS exploits this behavioral clustering through \textbf{Behavioral Archetype Factorization}.

Rather than training unique neural networks for each of the 1,000,000 agents, the offline distillation pipeline (detailed in Chapter \ref{ch:mdp_formulation}) trains and freezes compact, two-layer Multi-Layer Perceptrons (MLPs) representing four canonical market participant archetypes:

\begin{enumerate}
    \item \textbf{Market Makers (MM):} High-frequency liquidity providers operating two-sided quote ladders around the mid-price. Their policy balances spread capture against adverse selection risk ($\mu$) and inventory skew penalties ($\gamma \cdot (\text{Inventory}_t)^2$).
    \item \textbf{Institutional Whales (IW):} Capital-intensive liquidity consumers executing programmatic volume-weighted average price (VWAP) or percentage-of-volume (POV) liquidation schedules. Their policies balance market impact against execution urgency.
    \item \textbf{Momentum Scalpers (MS):} Latency-sensitive directional traders exploiting short-term Order Flow Imbalance (OFI) trends and volatility spikes.
    \item \textbf{Noise Traders (NT):} Retail and uncoordinated market participants characterized by stochastic liquidity shocks, loss aversion, and behavioral panic cascades governed by prospect theory \cite{kahneman1979prospect}.
\end{enumerate}

Each archetype policy $\pi_{\theta_k}(a | S_t)$ ($k \in \{\text{MM}, \text{IW}, \text{MS}, \text{NT}\}$) is parameterized by a lightweight 2-layer MLP with a 151-dimensional input layer, a 64-neuron hidden layer with fused ReLU activations, and a discrete action head outputting logits over the action tuple $A_t = (\text{action}_t, \Delta P_t, V_t)$. 

Across all four archetypes, the total parameter count remains under 5 million FP32 parameters ($<5\text{M}$ weights), translating to approximately 20 Megabytes ($\sim 20\text{ MB}$) of static weight data.

\begin{figure}[htbp]
\centering
\begin{tikzpicture}[
    scale=0.85, transform shape,
    box/.style={draw, rectangle, rounded corners, minimum width=2.8cm, minimum height=1.1cm, align=center},
    arr/.style={->, thick, >=stealth}
]
    \node[box, fill=blue!15] (state) {Agent State $S_t \in \mathbb{R}^{151}$\\(LOB, OFI, Inv, SIRC)};
    \node[box, fill=purple!15] (arch) [right=1.2cm of state] {Archetype Demux\\($k \in \{\text{MM, IW, MS, NT}\}$)\\(BMI2 Masking)};
    \node[box, fill=green!15] (mlp) [right=1.2cm of arch] {Shared 2-Layer MLP\\($<5\text{M}$ Params, $\sim 20\text{MB}$)\\(AVX2 / SSE4.2 FMA)};
    \node[box, fill=orange!15] (action) [right=1.2cm of mlp] {Action $A_t = (a, \Delta P, V)$\\(Sub-0.2 $\mu$s Latency)};

    \draw[arr] (state) -- (arch);
    \draw[arr] (arch) -- (mlp);
    \draw[arr] (mlp) -- (action);

    \node[draw, dashed, gray, fit=(arch) (mlp), label=below:\footnotesize Locked in Gracemont 4MB L2 Cluster Cache] (cachebox) {};
\end{tikzpicture}
\caption{Tier 2 distilled policy evaluation flow. Individual agent state vectors are evaluated against frozen archetype weights cached in the Gracemont L2 cache boundary.}
\label{fig:tier2_eval_flow}
\end{figure}

\section{State Buffer Layout and 10GB Memory Footprint}

While the neural network weights are globally shared across each archetype, individual agent instances require private, mutable execution state buffers to track financial positions, active orders, and historical metrics.

To eliminate heap fragmentation and garbage collection pauses, AMPS allocates a contiguous, data-oriented state buffer for all 1,000,000 agents in host DDR5 memory. Each individual agent instance is strictly constrained to an isolated state footprint of less than 10 Kilobytes ($<10\text{ KB}$):

\begin{equation}
    \text{Size}(\text{AgentState}) = \text{Portfolio} (64\text{ B}) + \text{OpenOrders} (256\text{ B}) + \text{LocalHistory} (512\text{ B}) + \text{PaddedBuffer} \le 10\text{ KB}
\end{equation}

Multiplying this per-agent buffer by 1,000,000 active instances yields a total system memory footprint of approximately 10.0 Gigabytes ($\sim 10.0\text{ GB}$) of DDR5 RAM. This allocation represents only 7.8\% of the 128GB physical memory available on the host workstation, leaving vast headroom for Zarr telemetry caches, `/dev/shm` shared memory rings, and the operating system kernel.

\section{Microstructural Justification for the Sub-0.2 Microsecond Target Latency}

In high-frequency quantitative simulations, policy evaluation speed is not merely a benchmark metric -- it is a prerequisite for microstructural determinism. 

At a target simulation throughput exceeding $> 100,000\text{ Ticks/Sec}$ across 1,000,000 active agents, the simulator processes tens of thousands of micro-decisions per millisecond. Limit Order Books enforce strict First-In, First-Out (FIFO) price-time priority \cite{ohara1995market, cartea2015algorithmic}. Every order submission or cancellation must be sequenced with microsecond-level causal ordering.

If an agent's policy evaluation latency $\tau_{\text{execution}}$ exceeds 200 nanoseconds ($0.2\,\mu\text{s}$), thread queues on host CPU processing pipelines experience stochastic queue delay skewing:

\begin{equation}
    \Delta t_{\text{queue}} = \sum_{i=1}^{N_{\text{pending}}} \tau_{\text{execution}}^{(i)} > \Delta t_{\text{clock\_tick}}
\end{equation}

When $\Delta t_{\text{queue}}$ exceeds the simulation tick interval, order arrivals arrive out of causal sequence, corrupting FIFO execution order and introducing artificial execution latency arbitrage. By bounding policy forward passes to $\tau_{\text{execution}} \le 0.2\,\mu\text{s}$, AMPS guarantees that all policy evaluations complete well within the scheduling quantum, preserving true price-time priority determinism across the matching engine.

\section{Gracemont E-Core Pinning and Cache Boundary Locking}

Executing 1,000,000 agents at sub-0.2 microsecond latency requires surgical hardware alignment. As established in Chapter \ref{ch:hardware_topography}, AMPS pins the Tier 2 execution engine strictly to the 12 Gracemont Efficiency Cores (Cores 16--27) of the Intel Core i7-14700K.

\subsection{L2 Cluster Locality}

The Gracemont architecture organizes E-cores into 4-core clusters, with each cluster sharing a dedicated 4MB Level 2 (L2) cache \cite{intel2026gracemont}. The shared archetype weights ($\sim 20\text{ MB}$ total, or $\sim 5\text{ MB}$ per archetype) and the Numba Limit Order Book data structures are statically pinned into these local 4MB L2 caches. 

Because the shared weights fit directly within the L2 cache boundaries, Gracemont E-cores never issue Level 3 (L3) cache fetches during standard policy evaluation forward passes. This architectural boundary completely insulates the high-frequency matching loop from L3 cache thrashing caused by Tier 1 LLM matrix multiplications executing on the Raptor Cove P-cores.

\subsection{Vectorized SIMD \& BMI2 Unpacking}

Forward-pass matrix-vector multiplications are executed via AVX2 and SSE4.2 Fused Multiply-Add (FMA) vector pipelines compiled directly through Numba `@njit(fastmath=True)`.

To unpack the 128-bit SIRC narrative signatures ($B_t \in \mathbb{R}^{128}$) from their compact 16-byte bitfields into the 151-dimensional state vector $S_t$, the engine utilizes Bit Manipulation Instruction Set 2 (BMI2) hardware intrinsics, specifically Parallel Extract (`PEXT`) and Parallel Deposit (`PDEP`):

\begin{lstlisting}[language=C, caption={BMI2 Bitwise Unpacking of 128-bit SIRC Signatures}]
// Direct hardware unpacking of 64-bit SIRC word using BMI2 PEXT
uint64_t mask = 0x5555555555555555ULL;
uint64_t unpacked_bits = _pext_u64(sirc_signature_high, mask);
\end{lstlisting}

These hardware intrinsics extract and deposit arbitrary bit patterns in a single CPU cycle, completely eliminating conditional branching (`if/else`) and bit-shift loops from the hot execution path.

Combined with AVX2 SIMD matrix-vector kernels, the Gracemont E-cores evaluate 2-layer MLP forward passes in an average of 142 nanoseconds per agent, comfortably satisfying the sub-0.2 microsecond latency invariant.
